% GENERATED FILE. Edit canonical lessons, then run npm run generate:books.
% canonical-source-hash: fnv1a64:836cbacb2fa9f60f
% canonical-lessons: JA-C159-kaji, JA-C159-suteru, JA-C159-wakeru, JA-C159-kasaneru, JA-C159-seiri

\chapter{Housework and Sorting the Rubbish}
\label{ch:ja-housework-and-sorting-the-rubbish}
\hlchaptermodality{159}

\begin{chapteropening}
\textbf{By the end of this chapter:} \emph{I can talk about housework, sorting and throwing things away.}

Complete the last lesson of chapter 159: I can talk about housework, sorting and throwing things away.
\end{chapteropening}

\section[\texorpdfstring{\ja{かじ} (kaji)}{かじ (kaji)}]{\ja{かじ} (kaji) — housework; chores}

\label{lesson:JA-C159-kaji}



\begin{quote}
Before the new one: say the Japanese for to open up, then the Japanese for dark.
\end{quote}

\subsection*{You'll want to know: \ja{かじ}}
\textbf{\ja{かじ}} — \emph{kaji} — \textquotedblleft{}housework; chores\textquotedblright{}.

\begin{quote}
\textbf{\ja{こどもも} \ja{かじを} \ja{てつだいます。}} — \emph{kodomo mo kaji o tetsudaimasu} — \textquotedblleft{}The children help with the housework too.\textquotedblright{}
\end{quote}

\begin{cousinweb}
\textbf{\ja{かじ}} (\emph{kaji}) is Sino-Japanese (kango): two Chinese characters read with their on'yomi, readings borrowed from Chinese. They are \emph{ka}, \textquotedblleft{}house, home\textquotedblright{} (Middle Chinese \emph{kae}) and \emph{ji}, \textquotedblleft{}matter, affair\textquotedblright{} (\emph{dzriH}) (KANJIDIC2; Middle Chinese in Baxter's transcription).
\end{cousinweb}

\begin{grammarlens}[title={What you've built}]
One more word for everyday life at A2.
\end{grammarlens}

\subsection*{Your turn}
\begin{itemize}
\raggedright
  \item \emph{Say it:} \emph{kaji}
  \item \emph{Say it:} \emph{kaji}, once more
  \item \emph{Say it:} \emph{kodomo mo kaji o tetsudaimasu} — The children help with the housework too.
  \item \emph{Recall:} say the Japanese for to open up, then the Japanese for dark, then say \emph{kaji} again
\end{itemize}

\begin{tcolorbox}[breakable,colback=teal!4,colframe=teal!35!black,title={Before you move on}]
What does \ja{かじ} mean? (\textquotedblleft{}Housework; chores\textquotedblright{}.) Say it once more.
\end{tcolorbox}
\section[\texorpdfstring{\ja{すてる} (suteru)}{すてる (suteru)}]{\ja{すてる} (suteru) — to throw away}

\label{lesson:JA-C159-suteru}



\begin{quote}
Before the new one: say the Japanese for dark, then the Japanese for housework.
\end{quote}

\subsection*{You'll want to know: \ja{すてる}}
\textbf{\ja{すてる}} — \emph{suteru} — \textquotedblleft{}to throw away\textquotedblright{}.

\begin{quote}
\textbf{\ja{まいしゅう} \ja{ごみを} \ja{すてます。}} — \emph{maishū gomi o sutemasu} — \textquotedblleft{}I throw out the rubbish every week.\textquotedblright{}
\end{quote}

Polite \textbf{\ja{すてます}}, \emph{sutemasu}; an ichidan verb.

\begin{cousinweb}
\textbf{\ja{すてる}} (\emph{suteru}) is a native Japanese word (wago), inherited rather than borrowed from Chinese; its kanji was attached later as a kun'yomi, and its older origin is not settled (JMdict, KANJIDIC2).
\end{cousinweb}

\begin{grammarlens}[title={What you've built}]
One more word for everyday life at A2.
\end{grammarlens}

\subsection*{Your turn}
\begin{itemize}
\raggedright
  \item \emph{Say it:} \emph{suteru}
  \item \emph{Say it:} \emph{suteru}, once more
  \item \emph{Say it:} \emph{maishū gomi o sutemasu} — I throw out the rubbish every week.
  \item \emph{Recall:} say the Japanese for dark, then the Japanese for housework, then say \emph{suteru} again
\end{itemize}

\begin{tcolorbox}[breakable,colback=teal!4,colframe=teal!35!black,title={Before you move on}]
What does \ja{すてる} mean? (\textquotedblleft{}To throw away\textquotedblright{}.) Say it once more.
\end{tcolorbox}
\section[\texorpdfstring{\ja{わける} (wakeru)}{わける (wakeru)}]{\ja{わける} (wakeru) — to sort; to divide}

\label{lesson:JA-C159-wakeru}



\begin{quote}
Before the new one: say the Japanese for housework, then the Japanese for to throw away.
\end{quote}

\subsection*{You'll want to know: \ja{わける}}
\textbf{\ja{わける}} — \emph{wakeru} — \textquotedblleft{}to sort; to divide\textquotedblright{}.

\begin{quote}
\textbf{\ja{ごみを} \ja{みっつに} \ja{わけます。}} — \emph{gomi o mittsu ni wakemasu} — \textquotedblleft{}I sort the rubbish into three kinds.\textquotedblright{}
\end{quote}

Polite \textbf{\ja{わけます}}, \emph{wakemasu}; an ichidan verb.

\begin{cousinweb}
\textbf{\ja{わける}} (\emph{wakeru}) is a native Japanese word (wago), inherited rather than borrowed from Chinese; its kanji was attached later as a kun'yomi, and its older origin is not settled (JMdict, KANJIDIC2).
\end{cousinweb}

\begin{grammarlens}[title={What you've built}]
One more word for everyday life at A2.
\end{grammarlens}

\subsection*{Your turn}
\begin{itemize}
\raggedright
  \item \emph{Say it:} \emph{wakeru}
  \item \emph{Say it:} \emph{wakeru}, once more
  \item \emph{Say it:} \emph{gomi o mittsu ni wakemasu} — I sort the rubbish into three kinds.
  \item \emph{Recall:} say the Japanese for housework, then the Japanese for to throw away, then say \emph{wakeru} again
\end{itemize}

\begin{tcolorbox}[breakable,colback=teal!4,colframe=teal!35!black,title={Before you move on}]
What does \ja{わける} mean? (\textquotedblleft{}To sort; to divide\textquotedblright{}.) Say it once more.
\end{tcolorbox}
\section[\texorpdfstring{\ja{かさねる} (kasaneru)}{かさねる (kasaneru)}]{\ja{かさねる} (kasaneru) — to stack; to pile one on another}

\label{lesson:JA-C159-kasaneru}



\begin{quote}
Before the new one: say the Japanese for to throw away, then the Japanese for to sort.
\end{quote}

\subsection*{You'll want to know: \ja{かさねる}}
\textbf{\ja{かさねる}} — \emph{kasaneru} — \textquotedblleft{}to stack; to pile one on another\textquotedblright{}.

\begin{quote}
\textbf{\ja{はこを} \ja{かさねます。}} — \emph{hako o kasanemasu} — \textquotedblleft{}I stack the boxes.\textquotedblright{}
\end{quote}

Polite \textbf{\ja{かさねます}}, \emph{kasanemasu}; an ichidan verb. Close to \textbf{\ja{つむ}}, \emph{tsumu}, \textquotedblleft{}to pile up\textquotedblright{}.

\begin{cousinweb}
\textbf{\ja{かさねる}} (\emph{kasaneru}) is a native Japanese word (wago), inherited rather than borrowed from Chinese; its kanji was attached later as a kun'yomi, and its older origin is not settled (JMdict, KANJIDIC2).
\end{cousinweb}

\begin{grammarlens}[title={What you've built}]
One more word for everyday life at A2.
\end{grammarlens}

\subsection*{Your turn}
\begin{itemize}
\raggedright
  \item \emph{Say it:} \emph{kasaneru}
  \item \emph{Say it:} \emph{kasaneru}, once more
  \item \emph{Say it:} \emph{hako o kasanemasu} — I stack the boxes.
  \item \emph{Recall:} say the Japanese for to throw away, then the Japanese for to sort, then say \emph{kasaneru} again
\end{itemize}

\begin{tcolorbox}[breakable,colback=teal!4,colframe=teal!35!black,title={Before you move on}]
What does \ja{かさねる} mean? (\textquotedblleft{}To stack; to pile one on another\textquotedblright{}.) Say it once more.
\end{tcolorbox}
\section[\texorpdfstring{\ja{せいり} (seiri)}{せいり (seiri)}]{\ja{せいり} (seiri) — sorting out; putting in order}

\label{lesson:JA-C159-seiri}



\begin{quote}
Before the new one: say the Japanese for to sort, then the Japanese for to stack.
\end{quote}

\subsection*{You'll want to know: \ja{せいり}}
\textbf{\ja{せいり}} — \emph{seiri} — \textquotedblleft{}sorting out; putting in order\textquotedblright{}.

\begin{quote}
\textbf{\ja{へやの} \ja{せいりは} \ja{たのしいです。}} — \emph{heya no seiri wa tanoshii desu} — \textquotedblleft{}Sorting out my room is fun.\textquotedblright{}
\end{quote}

\begin{cousinweb}
\textbf{\ja{せいり}} (\emph{seiri}) is Sino-Japanese (kango): two Chinese characters read with their on'yomi, readings borrowed from Chinese. They are \emph{sei}, \textquotedblleft{}put in order\textquotedblright{} (Middle Chinese \emph{tsyengX}) and \emph{ri}, \textquotedblleft{}reason, arrangement\textquotedblright{} (\emph{liX}) (KANJIDIC2; Middle Chinese in Baxter's transcription).
\end{cousinweb}

\begin{grammarlens}[title={What you've built}]
One more word for everyday life at A2.
\end{grammarlens}

\subsection*{Your turn}
\begin{itemize}
\raggedright
  \item \emph{Say it:} \emph{seiri}
  \item \emph{Say it:} \emph{seiri}, once more
  \item \emph{Say it:} \emph{heya no seiri wa tanoshii desu} — Sorting out my room is fun.
  \item \emph{Recall:} say the Japanese for to sort, then the Japanese for to stack, then say \emph{seiri} again
\end{itemize}

\begin{tcolorbox}[breakable,colback=teal!4,colframe=teal!35!black,title={Before you move on}]
What does \ja{せいり} mean? (\textquotedblleft{}Sorting out; putting in order\textquotedblright{}.) Say it once more.
\end{tcolorbox}
